% Couverture et quatrième de couverture du second livre : la mise en page du premier
% (docs/livre-methodes/couverture.tex), sur un aplat bleu de Prusse plutôt que cinabre,
% pour que les deux livres ne se confondent pas.
\definecolor{toncouv}{HTML}{1F4E6B}   % bleu de Prusse
\definecolor{toncouvencre}{HTML}{0C2232}

\newcommand{\couverture}{%
  \thispagestyle{empty}
  \begin{tikzpicture}[remember picture, overlay]
    \fill[toncouv] (current page.south west) rectangle (current page.north east);
    % des formules en filigrane, en trois ou quatre valeurs pour la profondeur
    \node[inner sep=0pt, text=toncouvencre, opacity=0.5, font=\Large, rotate=-4] at ($(current page.north west)+(3.4cm,-1.6cm)$) {$\displaystyle\sum_{n\geq 1}\frac{1}{n^2}=\frac{\pi^2}{6}$};
    \node[inner sep=0pt, text=white, opacity=0.18, font=\small, rotate=0] at ($(current page.north west)+(9.0cm,-1.4cm)$) {$\mathrm{Hom}(F\,X,\,Y)\cong\mathrm{Hom}(X,\,G\,Y)$};
    \node[inner sep=0pt, text=toncouvencre, opacity=0.3, font=\LARGE, rotate=3] at ($(current page.north west)+(8.6cm,-2.9cm)$) {$e^{i\pi}+1=0$};
    \node[inner sep=0pt, text=white, opacity=0.25, font=\small, rotate=2] at ($(current page.north west)+(2.9cm,-3.6cm)$) {$H^i(X,\mathcal F)=R^i\Gamma(X,\mathcal F)$};
    \node[inner sep=0pt, text=toncouvencre, opacity=0.22, font=\normalsize, rotate=-6] at ($(current page.north west)+(12.6cm,-2.2cm)$) {$\pi_1(S^1)\cong\mathbb Z$};
    \node[inner sep=0pt, text=toncouvencre, opacity=0.4, font=\footnotesize, rotate=0] at ($(current page.north west)+(5.6cm,-4.6cm)$) {\ttfamily example : (2:\!ℤ) \^{} 10 = 1024 := by norm\_num};
    \node[inner sep=0pt, text=white, opacity=0.15, font=\large, rotate=4] at ($(current page.north west)+(11.2cm,-4.3cm)$) {$\varprojlim_n \mathbb Z/p^n$};
    \node[inner sep=0pt, text=toncouvencre, opacity=0.45, font=\Large, rotate=3] at ($(current page.north west)+(4.0cm,-9.4cm)$) {$\displaystyle\int_{-\infty}^{+\infty} e^{-x^2}\,dx=\sqrt{\pi}$};
    \node[inner sep=0pt, text=white, opacity=0.22, font=\small, rotate=-3] at ($(current page.north west)+(10.0cm,-9.0cm)$) {$a^p\equiv a \pmod p$};
    \node[inner sep=0pt, text=toncouvencre, opacity=0.55, font=\footnotesize, rotate=0] at ($(current page.north west)+(6.6cm,-10.8cm)$) {\ttfamily theorem two\_mul (n : ℕ) : 2 * n = n + n := by omega};
    \node[inner sep=0pt, text=white, opacity=0.2, font=\normalsize, rotate=-5] at ($(current page.north west)+(2.8cm,-12.2cm)$) {$\det(AB)=\det A\,\det B$};
    \node[inner sep=0pt, text=toncouvencre, opacity=0.25, font=\LARGE, rotate=2] at ($(current page.north west)+(8.4cm,-12.6cm)$) {$\displaystyle\zeta(s)=\prod_{p}\frac{1}{1-p^{-s}}$};
    \node[inner sep=0pt, text=toncouvencre, opacity=0.5, font=\large, rotate=-3] at ($(current page.north west)+(4.2cm,-14.0cm)$) {$0\to A\to B\to C\to 0$};
    \node[inner sep=0pt, text=white, opacity=0.25, font=\small, rotate=5] at ($(current page.north west)+(10.4cm,-14.6cm)$) {$\nabla\cdot E=\rho/\varepsilon_0$};
    \node[inner sep=0pt, text=toncouvencre, opacity=0.35, font=\small, rotate=0] at ($(current page.north west)+(3.2cm,-15.8cm)$) {\ttfamily \#print axioms -- propext, Quot.sound};
    \node[inner sep=0pt, text=white, opacity=0.12, font=\Large, rotate=-2] at ($(current page.north west)+(7.4cm,-16.6cm)$) {$\chi(S^2)=V-E+F=2$};
    \node[inner sep=0pt, text=toncouvencre, opacity=0.28, font=\normalsize, rotate=4] at ($(current page.north west)+(2.9cm,-18.0cm)$) {$\displaystyle\binom{n}{k}=\binom{n-1}{k-1}+\binom{n-1}{k}$};
    \node[inner sep=0pt, text=toncouvencre, opacity=0.45, font=\footnotesize, rotate=-4] at ($(current page.north west)+(9.6cm,-18.2cm)$) {$\displaystyle\oint_\gamma \frac{dz}{z}=2\pi i$};
    % encore des formules, et le carré de fin de démonstration
    \node[inner sep=0pt, text=toncouvencre, opacity=0.35, font=\small, rotate=-2] at ($(current page.north west)+(11.4cm,-6.4cm)$) {$\mathcal O_X\to\mathcal F\to 0$};
    \node[inner sep=0pt, text=white, opacity=0.18, font=\footnotesize, rotate=3] at ($(current page.north west)+(6.4cm,-8.0cm)$) {$\Gamma(s)\Gamma(1-s)=\dfrac{\pi}{\sin\pi s}$};
    \node[inner sep=0pt, text=white, opacity=0.22, font=\normalsize, rotate=-4] at ($(current page.north west)+(11.6cm,-10.6cm)$) {$G/H\cong\mathrm{Gal}(L/K)$};
    \node[inner sep=0pt, text=white, opacity=0.2, font=\small, rotate=2] at ($(current page.north west)+(5.6cm,-19.6cm)$) {$\|x+y\|\le\|x\|+\|y\|$};
    \node[inner sep=0pt, text=toncouvencre, opacity=0.42, font=\large, rotate=-3] at ($(current page.north west)+(10.6cm,-16.0cm)$) {$\mathbb P^1(\mathbb C)\simeq S^2$};
    \node[inner sep=0pt, text=toncouvencre, opacity=0.28, font=\footnotesize, rotate=-2] at ($(current page.north west)+(2.6cm,-20.6cm)$) {$\forall\varepsilon>0,\ \exists\delta>0$};
    \node[inner sep=0pt, text=toncouvencre, opacity=0.22, font=\normalsize, rotate=4] at ($(current page.north west)+(11.4cm,-19.6cm)$) {$\mathrm{Spec}\,\mathbb Z$};
    \path[fill=toncouvencre, opacity=0.55] ($(current page.north west)+(12.0cm,-3.6cm)$) rectangle ++(3.2mm,3.2mm);
    \path[draw=white, line width=0.7pt, opacity=0.35] ($(current page.north west)+(7.0cm,-6.0cm)$) rectangle ++(2.2mm,2.2mm);
    \path[fill=white, opacity=0.25] ($(current page.north west)+(12.2cm,-8.0cm)$) rectangle ++(2.6mm,2.6mm);
    \path[draw=toncouvencre, line width=0.7pt, opacity=0.5] ($(current page.north west)+(3.2cm,-10.6cm)$) rectangle ++(4.0mm,4.0mm);
    \path[fill=toncouvencre, opacity=0.45] ($(current page.north west)+(11.6cm,-13.2cm)$) rectangle ++(2.4mm,2.4mm);
    \path[draw=white, line width=0.7pt, opacity=0.3] ($(current page.north west)+(8.6cm,-15.5cm)$) rectangle ++(3.0mm,3.0mm);
    \path[fill=white, opacity=0.22] ($(current page.north west)+(3.0cm,-16.9cm)$) rectangle ++(2.6mm,2.6mm);
    \path[draw=toncouvencre, line width=0.7pt, opacity=0.4] ($(current page.north west)+(12.0cm,-17.4cm)$) rectangle ++(3.6mm,3.6mm);
    \path[fill=toncouvencre, opacity=0.5] ($(current page.north west)+(9.0cm,-20.2cm)$) rectangle ++(2.8mm,2.8mm);
    \path[fill=white, opacity=0.3] ($(current page.north west)+(5.4cm,-2.6cm)$) rectangle ++(2.0mm,2.0mm);
    % le grand mot, couché, en encre sombre
    \node[rotate=90, anchor=south, inner sep=0pt, text=toncouvencre,
          font=\firathin\fontsize{118}{118}\selectfont]
      at ($(current page.south east)+(-0.9cm,12.4cm)$) {outillées};
    % le titre en petit, et le filet qui le relie au grand mot
    \node[anchor=base west, inner sep=0pt, text=white, font=\sffamily\fontsize{26}{30}\selectfont]
      (titre) at ($(current page.north west)+(2.2cm,-6cm)$) {Archives};
    \draw[toncouvencre, line width=1.6pt]
      ($(titre.base east)+(0.35cm,0.12cm)$) -- ($(current page.east|-titre.base)+(-4.6cm,0.12cm)$);
    \node[anchor=north west, inner sep=0pt, text=white!88, font=\sffamily\large,
          text width=10cm, align=left]
      at ($(titre.base west)+(0,-0.9cm)$) {reproduire la méthode sur des fonds publics};
    % les auteurs en bas à droite
    \node[anchor=south east, inner sep=0pt, text=white, font=\sffamily\normalsize,
          align=left]
      at ($(current page.south east)+(-5.4cm,2.2cm)$)
      {Michel Hua\\[2pt]{\small\color{white!80} avec Claude comme assistant}\\[2pt]
       {\footnotesize\color{white!70} version \versionlivre}};
    % la marque, discrète, en bas à gauche
    \node[anchor=south west, inner sep=0pt, text=white!75, font=\sffamily\scriptsize]
      at ($(current page.south west)+(2.2cm,2.2cm)$) {commutator.io};
  \end{tikzpicture}
  \null\clearpage
}

\newcommand{\quatrieme}{%
  \cleardoublepage
  \thispagestyle{empty}
  \begin{tikzpicture}[remember picture, overlay]
    \fill[toncouv] (current page.south west) rectangle (current page.north east);
    \node[anchor=north west, text width=11.6cm, align=justify, text=white,
          font=\normalsize, inner sep=0pt]
      at ($(current page.north west)+(2.4cm,-3.4cm)$) {%
      {\sffamily\Large Archives outillées}\par\bigskip
      Des bibliothèques et des universités ont mis en ligne, entiers et gratuits, les
      papiers de mathématiciens : brouillons, lettres, carnets, rédactions. Presque personne
      ne les lit. Un assistant de code peut les transcrire en LaTeX, à côté de leur
      fac-similé, et faire vérifier par Lean ce qu'ils démontrent.\par\medskip
      Ce livre décrit comment reprendre cette méthode sur un autre fonds. Il passe en revue,
      une à une, les archives sur lesquelles elle a servi : le fonds Grothendieck de
      Montpellier, les rédactions Bourbaki, les carnets de Quillen, les manuscrits de
      Gallica. Pour chacune, il dit comment on y accède, ce qui a coincé et ce qu'on a
      décidé. Il finit par un plan de travail pour les manuscrits mathématiques de
      Leibniz.\par\medskip
      Il complète \emph{Mathématiques outillées}, qui décrit les outils ; celui-ci décrit
      les sources.};
    \draw[toncouvencre, line width=1.6pt]
      ($(current page.south west)+(2.4cm,3.4cm)$) -- ++(3cm,0);
    \node[anchor=south west, inner sep=0pt, font=\sffamily\footnotesize, text=white!80]
      at ($(current page.south west)+(2.4cm,2.2cm)$) {version \versionlivre\ · 2026};
    \node[anchor=south east, inner sep=0pt, font=\sffamily\scriptsize, text=white!75]
      at ($(current page.south east)+(-2.4cm,2.2cm)$) {commutator.io};
  \end{tikzpicture}
  \null
}
